\chapter{Results and Discussion}
\label{chap:results}
% What was demonstrated and how well it worked. Only real measurements.
% Covers CRISP-DM MO.2, MO.4 and the Evaluation phase (EV.1-EV.3).

\section{Introduction}

\section{Experimental Setup and Test Design}
\crispdm{MO.2}
% MO.2 Generate test design, decided BEFORE assessing models.
% Exact environment (hardware, software versions, dates), data split
% (proportions, stratification, temporal/group split, leakage prevention),
% validation protocol, baseline, and how the test set was kept untouched.
% Split sizes must add up with the data description in Chapter 3.

\section{Evaluation Metrics}
% Each metric as a numbered equation with every variable defined, plus a
% plain-language explanation of what it rewards or penalizes.

\section{Model Assessment}
\crispdm{MO.4}
% MO.4 Assess model: results per split (validation vs held-out test), variants
% compared, tuning iterations, against the technical criteria of Chapter 1.
% Plots and tables generated from the real logs only.

\section{Functional Demonstration}
% Qualitative evidence: real screenshots, example predictions, representative
% successes and failures (anonymized).

\section{Evaluation Against Business Success Criteria}
\label{sec:business-evaluation}
\crispdm{EV.1}
% EV.1 Evaluate results: for EACH business success criterion of Chapter 1, is
% it met, not met, or not measurable yet, and on what evidence? Which model is
% approved and why? This is a business judgement; do not just repeat the
% metrics above.

\section{Review of the Process}
\crispdm{EV.2}
% EV.2 Review process: error analysis, steps skipped or rushed, leakage checks,
% risks from the situation assessment that materialized, threats to validity,
% limitations.

\section{Discussion and Next Steps}
\crispdm{EV.3}
% EV.3 Determine next steps: the possible actions (deploy, iterate, collect more
% data, stop), the decision taken and its reasons. Must agree with the
% deployment plan and the General Conclusion.

\section{Conclusion}
